\chapter{Wstęp}

We współczesnym świecie, zdominowanym przez koncepcję Internetu Rzeczy (\textit{ang. Iot - Internet of Things}), robotykę oraz pojazdy autonomiczne, transmisja danych na odległość stała się zagadnieniem powszechnym.
O ile klasyczne przesyłanie informacji, takie jak pobieranie plików, czy też przeglądanie aplikacji i stron internetowych, opiera się na gwarancji bezbłędności transmisji, niezależnie od czasu trwania tego procesu, o tyle telemetria i zdalne sterowanie wprowadzają radykalnie odmienny paradygmat.
W systemach czasu rzeczywistego (\textit{ang. Real-Time Systems}) wartość informacji drastycznie spada wraz z upływem czasu.

Na bezwzględny czas transmisji danych składa się wiele czynników.
Zjawiska takie jak opóźnienia w buforach, zmienność opóźnień, czy utrata ciągłości sterowania potrafią doprowadzić do katastrofalnych skutków, na przykład utraty kontroli nad lecącym bezzałogowym statkiem powietrznym.
Inżynierowie projektujący tego typu rozwiązania nieustannie stają przed dylematem wyboru odpowiedniego medium transmisyjnego oraz protokołu warstwy transportowej lub aplikacyjnej.
Gwarancja dostarczenia każdego pakietu zazwyczaj wiąże się z niedopuszczalnym w systemach dynamicznych narzutem czasowym, podczas gdy minimalizacja opóźnień wymaga akceptacji nieuchronnego gubienia części danych w eterze.
